\documentclass[11pt]{article}
\usepackage{amsmath,amssymb,amsthm,amsfonts,hyperref,booktabs}
\newtheorem{theorem}{Theorem}[section]
\newtheorem{proposition}[theorem]{Proposition}
\newtheorem{lemma}[theorem]{Lemma}
\newtheorem{corollary}[theorem]{Corollary}
\newtheorem{definition}[theorem]{Definition}
\newtheorem{remark}[theorem]{Remark}
\newtheorem{observation}[theorem]{Observation}
\providecommand{\C}{\mathbb{C}}
\providecommand{\Z}{\mathbb{Z}}
\providecommand{\cA}{\mathcal{A}}
\providecommand{\cO}{\mathcal{O}}
\providecommand{\Coh}{\mathrm{Coh}}
\providecommand{\HH}{\mathrm{HH}}
\providecommand{\HC}{\mathrm{HC}}
\providecommand{\HKR}{\mathrm{HKR}}
\providecommand{\Tot}{\mathrm{Tot}}
\providecommand{\Hilb}{\mathrm{Hilb}}
\providecommand{\Hom}{\mathrm{Hom}}
\providecommand{\rk}{\mathrm{rk}}
\providecommand{\Ker}{\mathrm{Ker}}
\providecommand{\str}{\mathrm{str}}
\providecommand{\vol}{\mathrm{vol}}
\title{Dimensional Stratification of $\kappa_{\mathrm{ch}}$ on Compact Calabi--Yau Manifolds}
\author{Raeez Lorgat}
\date{2026-04-22}
\begin{document}
\maketitle

\begin{abstract}
The functor $\Phi_d\colon \mathrm{CY}_d\text{-}\mathrm{cat}\to\mathrm{ChirAlg}$
attaches to a compact Calabi--Yau $d$-fold $X$ a chiral algebra $\cA_X$
whose degree-two shadow scalar $\kappa_{\mathrm{ch}}(\cA_X)$ is a
first-principles invariant of the $(0,\bullet)$ Hodge column.
The scalar is the Hodge-filtered supertrace
$\Xi(X)=\sum_{q=0}^{d}(-1)^{q}h^{0,q}(X)$. On this identification the
Serre involution $\sigma\colon q\mapsto d-q$, with character $(-1)^d$
on alternating sums, forces $\Xi(X)=0$ at every odd $d$ and partitions
the even-$d$ landscape into the strict-Calabi--Yau stratum ($\Xi=2$)
and the irreducible holomorphic-symplectic stratum ($\Xi=(d/2)+1$).
The identity $\kappa_{\mathrm{ch}}=\chi(\cO_X)$ is a structural equality
at K3 alone, a tautology of Serre cancellation at every other odd-$d$
compact case, and a coincidence of middle columns at $K3^{[n]}$. The
correct stratified formula separates compact ($\Xi$ on $\HC^{-}_{d}$),
local-surface ($\tfrac12\chi_{\mathrm{top}}$ on compactly supported
Dolbeault), and conifold (rank of the compact generator on
Kronecker-quiver Hochschild) regimes. A six-dimensional holomorphic
Chern--Simons avatar realises $\kappa_{\mathrm{ch}}$ as a Costello
one-loop BV--BRST trace on the ghost Dolbeault complex of $X$; odd-$d$
vanishing is a fermionic Serre symmetry of the ghost sector, and the
boundary $E_1$-chiral algebra on $C$ is the $E_3$-Koszul averaging of
the $E_3$-algebra structure on compactly supported Dolbeault of the
$6$-fold $Z=X\times C$.
\end{abstract}

\section{Four scalars, one functor}
\label{sec:four-scalars}

The CY-to-chiral functor factorises through a holomorphic factorisation
stage,
\[
  \Phi_d \;=\; \mathrm{Sp}^{\mathrm{ch}}_{\Sigma_{d-1},C}\circ\Phi^{\mathrm{FA}}_d
  \colon\;\mathrm{CY}_d\text{-}\mathrm{cat}
  \;\longrightarrow\;E_d\text{-}\mathrm{HolFA}(X)
  \;\longrightarrow\;E_1\text{-}\mathrm{HolFA}(C).
\]
Four scalars attach to the pair $(X,\cA_X)$, each native to a different
stage and each obeying a different Grothendieck--Riemann--Roch-style
rule on products.
\begin{itemize}
\item $\kappa_{\mathrm{ch}}(\cA_X)$, the degree-two shadow coefficient
      of the factorisation algebra $\Phi^{\mathrm{FA}}_d(D^{b}(\Coh(X)))$;
      native to the $\Phi^{\mathrm{FA}}_d$ stage. K\"unneth-multiplicative
      on compact products: $\Xi(X\times Y)=\Xi(X)\cdot\Xi(Y)$.
\item $\kappa_{\mathrm{cat}}(X)=\chi(\cO_X)=\sum_{q}(-1)^{q}h^{0,q}(X)$,
      the holomorphic Euler characteristic of the variety, also
      K\"unneth-multiplicative:
      $\chi(\cO_{X\times Y})=\chi(\cO_X)\cdot\chi(\cO_Y)$.
      Equal to $\Xi(X)$ as a number; the two invariants differ in
      construction (one natively on $\mathrm{HC}^{-}_{d}$, the other on
      the variety) and in the role of Serre duality (a tautology for
      $\chi$, a character-bearing involution for $\Xi$).
\item $\kappa_{\mathrm{BKM}}(\Phi_N)=c_N(0)/2$, the weight of the
      Borcherds denominator product attached to a K3-fibred CY$_3$;
      native to the $\mathrm{Sp}^{\mathrm{ch}}$ specialisation
      (Gritsenko 1999; Borcherds 1995, 1998).
\item $\kappa_{\mathrm{fiber}}$, the rank (or Euler characteristic)
      of the fibre of the specialisation cycle.
\end{itemize}

On $K3\times E$, K\"unneth yields
\[
   \chi(\cO_{K3\times E})\;=\;\chi(\cO_{K3})\cdot\chi(\cO_E)\;=\;2\cdot 0\;=\;0,
\]
a scalar of the total space --- not the K3-fibre value
$\chi(\cO_{K3})=2$. The two must never be conflated. In the
Hodge-supertrace lane $\Xi(K3\times E)=(1-1+1-1)=0$ as well, coinciding
numerically with $\chi(\cO_{K3\times E})$ but via a distinct
structural mechanism: an uncancelled Serre pair at $q\in\{1,2\}$ against
two units at $q\in\{0,3\}$, not the trivial $\chi=0$ of a $q$-constant
column.

The additive generator-count scalar
$\kappa_{\mathrm{ch}}^{\mathrm{Heis}}(\cA_X)$, attached to the
Heisenberg level on categorical $\HH^{\mathrm{cat},\bullet+1}$, is a
distinct invariant again: on $K3\times E$ it takes the value
$\kappa_{\mathrm{ch}}^{\mathrm{Heis}}(K3)+\kappa_{\mathrm{ch}}^{\mathrm{Heis}}(E)=2+1=3$
by additive generator counting. It is neither the Hodge supertrace nor
the Euler characteristic; it is the count of independent Heisenberg
bosons on the $\Phi^{\mathrm{FA}}_d$-output lattice.

The five scalars $\{0,3,5,24\}$ of the K3 landscape are the values of
four distinct constructions (Hodge supertrace, Heisenberg generator
count, Borcherds weight, fibre rank) on one compact Calabi--Yau
threefold. The distinct values come from distinct constructions, not
from four applications of $\Phi$ to one object.

\section{The Hodge supertrace}
\label{sec:hodge-supertrace}

Let $X$ be a compact Calabi--Yau manifold of complex dimension $d$ with
$\omega_X\simeq\cO_X$. Write $h^{p,q}=h^{p,q}(X)=\dim_{\C}H^q(X,\Omega^p_X)$
for the Hodge numbers. The Serre functor on $D^{b}(\Coh(X))$ is the
shift $[d]$; restricted to the structure-sheaf column this reads
$h^{0,q}=h^{0,d-q}$.

\begin{definition}[Hodge supertrace]
\label{def:xi}
The Hodge-filtered supertrace of $X$ is
\[
   \Xi(X)\;:=\;\sum_{q=0}^{d}(-1)^{q}\,h^{0,q}(X).
\]
\end{definition}

\begin{theorem}[Supertrace identification]
\label{thm:supertrace}
Let $\cA_X=\Phi_d(D^{b}(\Coh(X)))$ be the chiral algebra of the
CY-to-chiral correspondence. Then
\[
   \kappa_{\mathrm{ch}}(\cA_X)\;=\;\Xi(X),
\]
at generic parameters of $\Phi_d$.
\end{theorem}

\begin{proof}
By Hochschild--Kostant--Rosenberg, the categorical Hochschild
homology decomposes as
$\HH_{\bullet}(D^{b}(\Coh(X)))\simeq
\bigoplus_{p+q=\bullet}H^{q}(X,\Omega^{p}_{X})$,
with the Connes $B$-operator intertwined with the de Rham differential
up to Caldararu twist. The Calabi--Yau structure determines the negative
cyclic lift $\HC^{-}_{d}$, the $S^{d}$-framing datum of $\Phi_d$. The
degree-two shadow $S_{2}(\cA_X)$ is the value of the chiral $r$-matrix
pairing on $F^{0}/F^{-1}\cong\bigoplus_{q}H^{q}(X,\cO_X)$, the
structure-sheaf column; the Mukai $r$-matrix (Proposition
\ref{prop:serre-duality-cy}-analog) descends to the identity on this
column. Each summand $H^{q}(X,\cO_X)$ enters with cohomological sign
$(-1)^{q}$ from the chiral bar complex. Summing,
\[
   \kappa_{\mathrm{ch}}(\cA_X)\;=\;\sum_{q=0}^{d}(-1)^{q}\dim H^{q}(X,\cO_X)
   \;=\;\sum_{q=0}^{d}(-1)^{q}h^{0,q}(X)\;=\;\Xi(X). \qedhere
\]
\end{proof}

\begin{lemma}[Character of the Serre involution on alternating sums]
\label{lem:serre-sign}
The involution $\sigma\colon q\mapsto d-q$ on $\{0,1,\dots,d\}$ acts on
the alternating sum $\Xi(X)$ by the scalar $(-1)^{d}$:
\[
   \sigma\cdot\Xi(X)\;=\;\sum_{q}(-1)^{d-q}h^{0,d-q}
   \;=\;(-1)^{d}\sum_{q}(-1)^{q}h^{0,q}
   \;=\;(-1)^{d}\Xi(X),
\]
using Serre duality $h^{0,q}=h^{0,d-q}$. At odd $d$ the character is
$-1$, forcing $\Xi(X)=-\Xi(X)$ and hence $\Xi(X)=0$. At even $d$ the
character is $+1$ and the identity is tautological; $\Xi(X)$ is then
controlled by the middle columns.
\end{lemma}

\begin{proof}
Substitute $q\leftrightarrow d-q$ in Definition~\ref{def:xi}. Serre
duality identifies $h^{0,d-q}=h^{0,q}$. At $d$ odd the involution has no
fixed point (the would-be fixed value $q=d/2$ is non-integral); every
pair $\{q,d-q\}$ contributes $(-1)^{q}+(-1)^{d-q}=(-1)^{q}(1+(-1)^{d})=0$.
At $d$ even, pairs with $q\ne d/2$ contribute
$(-1)^{q}\cdot 2h^{0,q}$ and the unique fixed point $q=d/2$ contributes
$(-1)^{d/2}h^{0,d/2}$.
\end{proof}

Lemma~\ref{lem:serre-sign} is the root cause of the dimensional
stratification. Not ``additive under direct sums but not fibre
products'' --- the Hodge supertrace is K\"unneth-\emph{multiplicative},
as is the holomorphic Euler characteristic, and the two numerically
coincide on compact $X$. The content of the stratification is a single
$\Z/2$-symmetry: Serre duality as a fermionic involution on the
$(0,\bullet)$ column, with character $(-1)^d$ on alternating traces.

\section{Stratification}
\label{sec:stratification}

\begin{theorem}[Calabi--Yau tri-stratum]
\label{thm:tri-stratum}
Let $X$ be a compact CY$_d$ with $h^{1,0}(X)=0$. One of three structural
strata holds, determined by the parity of $d$ and the
Beauville--Bogomolov holonomy type.
\begin{enumerate}
\item[\textup{(I)}] \emph{Odd-$d$ stratum.} $\Xi(X)=0$ for every compact
      $X$. Mechanism: fixed-point freeness of $\sigma$
      (Lemma~\ref{lem:serre-sign}).
\item[\textup{(II)}] \emph{Strict Calabi--Yau stratum (holonomy
      $\mathrm{SU}(d)$, $d$ even).} $h^{0,k}=0$ for $0<k<d$ and
      $\Xi(X)=2$.
\item[\textup{(III)}] \emph{Irreducible holomorphic-symplectic
      stratum (holonomy $\mathrm{Sp}(n)$, $d=2n$).} $h^{0,2k}=1$ for
      $0\le k\le n$, $h^{0,q}=0$ for odd $q$, and
      $\Xi(X)=n+1=(d/2)+1$.
\end{enumerate}
Strata \textup{(II)} and \textup{(III)} coincide at $d=2$
($\mathrm{SU}(2)=\mathrm{Sp}(1)$); for $d\ge 4$ they are disjoint.
Together with \textup{(I)}, they exhaust the Beauville--Bogomolov
classification on connected simply-connected $X$ with $h^{1,0}=0$.
\end{theorem}

\begin{proof}
Stratum (I) is Lemma~\ref{lem:serre-sign}. In stratum (II), the only
nonzero $h^{0,q}$ are $q=0$ and $q=d$, each equal to $1$; $\Xi(X)=1+(-1)^{d}=2$
for $d$ even. In stratum (III), $\Xi(X)=\sum_{k=0}^{n}h^{0,2k}=n+1$.
Mutual exclusivity uses the Beauville--Bogomolov decomposition of
compact K\"ahler CY manifolds with $c_{1}=0$ into torus, strict-CY, and
holomorphic-symplectic factors (Beauville 1983; Bogomolov 1974); the
hypothesis $h^{1,0}=0$ eliminates the torus factor.
\end{proof}

\begin{corollary}[Running values]
\label{cor:running-values}
The Hodge supertrace takes the following values on the standard compact
Calabi--Yau examples.
\begin{center}
\begin{tabular}{lcc}
\toprule
$X$ & $(d,h^{0,\bullet})$ & $\Xi(X)=\kappa_{\mathrm{ch}}(\cA_X)$\\
\midrule
Elliptic curve $E$ & $(1,(1,1))$ & $0$\\
K3 & $(2,(1,0,1))$ & $2$\\
Abelian surface & $(2,(1,2,1))$ & $0$\\
Bielliptic surface & $(2,(1,1,0))$ & $0$\\
Quintic $X_{5}\subset\mathbb{P}^{4}$ & $(3,(1,0,0,1))$ & $0$\\
$K3\times E$ & $(3,(1,1,1,1))$ & $0$\\
$E^{3}$ & $(3,(1,3,3,1))$ & $0$\\
Sextic $X_{6}\subset\mathbb{P}^{5}$ & $(4,(1,0,0,0,1))$ & $2$\\
$K3^{[2]}$ & $(4,(1,0,1,0,1))$ & $3$\\
Septic $X_{7}\subset\mathbb{P}^{6}$ & $(5,(1,0,0,0,0,1))$ & $0$\\
$K3\times X_{5}$ & $(5,(1,0,1,1,0,1))$ & $0$\\
$K3\times K3$ & $(4,(1,0,2,0,1))$ & $4$\\
$K3\times K3^{[2]}$ & $(6,(1,0,2,0,2,0,1))$ & $6$\\
\bottomrule
\end{tabular}
\end{center}
\end{corollary}

The identification $\Xi(X)=\chi(\cO_X)$ as a \emph{number} is a
tautology. The identification as a \emph{structural equality with
content} --- both sides counting nonzero entries of a single Hodge
column with compatible signs, neither forced to zero by Serre
cancellation --- holds on exactly one row: K3. Every odd-$d$ row is
stratum-(I) Serre cancellation, not a content equality. At $K3^{[2]}$
the column $(1,0,1,0,1)$ sums to $\Xi=3$, matching
$\chi(\cO_{K3^{[2]}})=3$ (Beauville--G\"ottsche); but the mechanism
on the $\Phi$-side is three surviving middle-column contributions in
stratum (III), not a two-entry match.

K\"unneth multiplicativity $\Xi(X\times Y)=\Xi(X)\cdot\Xi(Y)$ is the
value at $t=-1$ of the multiplicative Hodge--Poincar\'e polynomial of
the $(0,\bullet)$ column; it coincides with the K\"unneth
multiplicativity of $\chi(\cO_{\bullet})$ on total spaces. The two
rows $K3\times K3$ and $K3\times K3^{[2]}$ in the table illustrate the
phenomenon: $\Xi(K3\times K3)=2\cdot 2=4=\chi(\cO_{K3\times K3})$ and
$\Xi(K3\times K3^{[2]})=2\cdot 3=6=\chi(\cO_{K3\times K3^{[2]}})$; the
multiplicative rule is a theorem for both sides, not a hypothesis.

\section{Odd-$d$: three falsifications of $\kappa_{\mathrm{ch}}=\chi(\cO_X)$}
\label{sec:counterexamples}

At $d=3$ the identity $\kappa_{\mathrm{ch}}=\chi(\cO_X)$ is a tautology
on every compact CY$_3$ with $h^{1,0}=0$ (both sides are zero by
Lemma~\ref{lem:serre-sign}). The interesting falsifications are
non-compact toric geometries on which the first-stage scalar carries
real content.

\begin{proposition}[Local $\mathbb{P}^2$]
\label{prop:local-p2}
Let $X=\Tot(\cO_{\mathbb{P}^2}(-3)\to\mathbb{P}^2)$, the canonical bundle
of $\mathbb{P}^2$. Then
\[
   \kappa_{\mathrm{ch}}(\cA_X)\;=\;\tfrac{3}{2}.
\]
\end{proposition}

\begin{proof}
The base $S=\mathbb{P}^2$ is Fano; $X=\Tot(\omega_S\to S)$ is a
non-compact local CY$_3$. Beilinson's tilting resolution of $\mathbb{P}^2$
produces the generator
$T=\cO\oplus\cO(1)\oplus\cO(2)\in D^{b}(\Coh(\mathbb{P}^2))$ with
$\Hom(T,T)=\mathrm{End}(T)$ the path algebra of the Beilinson quiver
$Q_{\mathbb{P}^2}$ (three vertices, three arrows
$\cO\xrightarrow{x_{0},x_{1},x_{2}}\cO(1)\xrightarrow{y_{0},y_{1},y_{2}}\cO(2)$,
three arrows
$\cO\to\cO(2)$ through the composition). On $X=\Tot(\omega_S)$ the
tilting generator pulls back to a Kontsevich--Soibelman
quiver-with-potential $(\widetilde{Q},W)$ where $\widetilde{Q}$ is the
$\Z/3$-cyclic double of $Q_{\mathbb{P}^2}$ and $W$ is the superpotential
$W=\sum\epsilon_{ijk}x_{i}y_{j}z_{k}$ realising the $\Z/3$-McKay
geometry at the origin (Klebanov--Witten 1998; Ooguri--Yamazaki 2008).

The degree-two shadow scalar on $\cA_X=\Phi_3(D^{b}(\Coh(X)))$ reduces,
by the Kontsevich--Soibelman CoHA identification, to a supertrace on
compactly supported Dolbeault cohomology of $X$. Serre pairing between
compact and non-compact Dolbeault columns of
$X=\Tot(\omega_S)$ exchanges $H^{q}_{c}(X,\cO_X)$ with
$H^{d-q}(X,\omega_X)^{\vee}=H^{3-q}(X,\cO_X)^{\vee}$; on $\Tot(\omega_S)$
this collapses to the base-level supertrace
\[
   \Xi^{c}(X)\;=\;\tfrac{1}{2}\chi_{\mathrm{top}}(S)\;=\;\tfrac{1}{2}\cdot 3\;=\;\tfrac{3}{2},
\]
where the factor $\tfrac12$ arises because the fibre direction
contributes a half-Euler pairing: compact fibres $\omega_S\to\mathrm{pt}$
reduce Dolbeault on $X$ to a fermion-halved copy of Dolbeault on $S$
(the standard local-CY$_3$ reduction of Szendr\H{o}i 2008,
Nagao--Nakajima 2011).

Independent confirmation: direct McKay computation on the
$\Z/3$-orbifold point of the Kontsevich--Soibelman CoHA yields
$\kappa_{\mathrm{ch}}=3/2$ via the three half-modes of the rank-$3$
lattice VOA associated to the Beilinson basis. The $\Z/3$-McKay
potential forces an infinite shadow tower (class M); the degree-two
scalar survives at $3/2$ independently of the tower.
\end{proof}

\begin{proposition}[Resolved conifold]
\label{prop:conifold}
Let $X_{\mathrm{con}}=\Tot(\cO(-1)^{\oplus 2}\to\mathbb{P}^1)$ be the
resolved conifold. Then
\[
   \kappa_{\mathrm{ch}}(\cA_{X_{\mathrm{con}}})\;=\;1.
\]
\end{proposition}

\begin{proof}
The Kontsevich--Soibelman quiver-with-potential for $X_{\mathrm{con}}$
is the Kronecker--$A_{1}$ quiver: two vertices $\{\bullet_{1},\bullet_{2}\}$,
two pairs of anti-parallel arrows $A_{1,2}\colon\bullet_{1}\rightleftarrows\bullet_{2}$
and $B_{1,2}\colon\bullet_{1}\rightleftarrows\bullet_{2}$, and zero
superpotential $W=0$ (Szendr\H{o}i 2008, Nagao 2011). The category
$D^{b}(\Coh(X_{\mathrm{con}}))$ admits a single compact generator
$\cO_{\mathbb{P}^{1}}$: the exceptional $\mathbb{P}^{1}$ is the unique
compact cycle in the small resolution. The categorical Hochschild
complex is that of modules over the Kronecker quiver path algebra with
$W=0$; its degree-zero part is $\C^{2}$ (idempotents at the two
vertices) and its Mukai pairing is the standard rank-$1$ hyperbolic
lattice.

The shadow tower produces a rank-$1$ Heisenberg VOA $\mathcal{H}_{1}$
at level $1$: a single free-boson mode $\alpha(z)$ with
$\alpha(z)\alpha(w)\sim 1/(z-w)^{2}$, attached to the unique compact
generator. The degree-two shadow $\kappa_{\mathrm{ch}}$ of a
Heisenberg VOA $\mathcal{H}_{k}$ equals its level $k$; at
$X_{\mathrm{con}}$ the level is $1$.
\end{proof}

\begin{proposition}[Generic quintic]
\label{prop:quintic}
Let $X_{5}\subset\mathbb{P}^{4}$ be a generic smooth quintic threefold.
Two invariants share the name $\kappa_{\mathrm{ch}}$; they attach to
different constructions and take different values.
\begin{itemize}
\item Hodge-supertrace lane:
      $\kappa_{\mathrm{ch}}^{\mathrm{Hodge}}(\cA_{X_{5}})=\Xi(X_{5})=0$,
      via the Serre cancellation $(1,0,0,1)$ of
      Lemma~\ref{lem:serre-sign}.
\item Bershadsky--Cecotti--Ooguri--Vafa (BCOV) lane:
      $\kappa_{\mathrm{ch}}^{\mathrm{BCOV}}(\cA_{X_{5}})=\chi(X_{5})/24=-200/24=-25/3$,
      the one-loop free energy of the topological string on $X_{5}$
      (Bershadsky--Cecotti--Ooguri--Vafa 1994; Klemm--Pandharipande 2007).
\end{itemize}
Neither equals $\chi(\cO_{X_{5}})=0$ as a structural identity: the first
equals zero by Serre cancellation, the second equals $-25/3$ by a
one-loop formula unrelated to the structure sheaf.
\end{proposition}

\begin{proof}
The Hodge column $(1,0,0,1)$ is Serre-paired, so
$\Xi(X_{5})=1-0+0-1=0$ by Lemma~\ref{lem:serre-sign}. The BCOV one-loop
free energy on a compact CY$_{3}$ is
$F_{1}(X)=\chi(X)/24$, a moduli-independent constant on the
Bogomolov--Tian--Todorov moduli. For the quintic
$\chi(X_{5})=\int_{X_{5}}c_{3}(TX_{5})=-200$, giving $F_{1}=-25/3$; this
is $\kappa_{\mathrm{ch}}^{\mathrm{BCOV}}$ at the Gepner point.
\end{proof}

Three odd-$d$ geometries falsify three distinct naive guesses. Local
$\mathbb{P}^2$ produces the non-integral scalar $3/2$, which cannot
equal any compact Euler characteristic. The resolved conifold produces
$1$, a value assigned not by a Hodge column at all (the conifold base
is a curve, not a surface) but by the rank of the single compact
generator of the Kronecker Hochschild complex. The compact quintic
supports two legitimate scalars, neither equal to $\chi(\cO_{X_{5}})$
outside the trivial Serre sense.

\section{The correct stratified formula}
\label{sec:correct-formula}

\begin{theorem}[Stratified formula for $\kappa_{\mathrm{ch}}$]
\label{thm:stratified}
Let $X$ be a Calabi--Yau variety of complex dimension $d$ with
$h^{1,0}=0$ on the compact part. Then
\[
   \kappa_{\mathrm{ch}}(\cA_X)\;=\;\begin{cases}
   \displaystyle\sum_{q=0}^{d}(-1)^{q}h^{0,q}(X) & X\text{ compact},\\[4pt]
   \tfrac{1}{2}\chi_{\mathrm{top}}(S) & X=\Tot(\omega_{S}\to S),\ S\text{ Fano surface},\\[4pt]
   \rk(\text{compact generator})=1 & X=X_{\mathrm{con}}\text{ resolved conifold}.
   \end{cases}
\]
In each case the right-hand side is a supertrace computed on a different
chain complex: the negative cyclic complex $\HC^{-}_{d}$ in the compact
case; compactly supported Dolbeault in the local-surface case; the
Hochschild homology of the Kronecker-$A_{1}$ quiver path algebra in the
conifold case.
\end{theorem}

\begin{proof}
The compact case is Theorem~\ref{thm:supertrace}. The local-surface
case replaces the compact Serre pairing on Dolbeault cohomology with
compactly supported Serre on $X=\Tot(\omega_{S})$; this reduces the
$(0,\bullet)$-supertrace on $X$ to a half-Euler characteristic of the
base $S$, the factor $\tfrac12$ coming from the fermionic fibre
direction (Proposition~\ref{prop:local-p2}; Nagao--Nakajima 2011). The
conifold case is Proposition~\ref{prop:conifold}.
\end{proof}

\begin{corollary}[When $\kappa_{\mathrm{ch}}=\chi(\cO_X)$ has content]
\label{cor:identity-scope}
The identity $\kappa_{\mathrm{ch}}(\cA_X)=\chi(\cO_X)$ with both sides
nonzero and no Serre cancellation on either side occurs, within compact
Calabi--Yau manifolds with $h^{1,0}=0$, exactly at $(d,h^{1,0},\text{Hodge column})=(2,0,(1,0,1))$,
i.e.\ at K3. On every irreducible holomorphic-symplectic manifold of
dimension $d=2n\ge 4$, both sides equal $n+1>2$ but by the middle-column
mechanism (Beauville--G\"ottsche) on the $\chi$-side and the surviving
three-Hodge-contribution mechanism (stratum (III)) on the
$\Xi$-side. At every compact odd-$d$ example, the identity is a
tautology of Serre cancellation. At local $\mathbb{P}^2$ and the
resolved conifold the identity fails nontrivially (fractional and
compact-generator values, respectively).
\end{corollary}

\begin{observation}[What the root cause actually is]
\label{obs:root-cause}
The stratification obstruction is a single $\Z/2$-symmetry: the Serre
involution $\sigma\colon q\mapsto d-q$, acting on alternating sums with
character $(-1)^{d}$, forces $\Xi(X)=0$ at odd $d$. K\"unneth
multiplicativity, not additivity, governs the compact-product calculus:
$\Xi(X\times Y)=\Xi(X)\cdot\Xi(Y)$, inherited from
$h^{0,\bullet}(X\times Y)=h^{0,\bullet}(X)\boxtimes h^{0,\bullet}(Y)$
and evaluation at $t=-1$ of the bivariate Hodge--Poincar\'e polynomial.
The appearance of an additive rule comes from a different scalar ---
the Heisenberg-level generator count
$\kappa_{\mathrm{ch}}^{\mathrm{Heis}}$ on the categorical
$\HH^{\mathrm{cat},\bullet+1}$-stage --- and is native to that
construction, not to the Hodge supertrace. The additive reading
attached in the K3 landscape to $K3\times E$
($\kappa_{\mathrm{ch}}^{\mathrm{Heis}}(K3\times E)=2+1=3$) and the
multiplicative reading of the Hodge supertrace
($\Xi(K3\times E)=2\cdot 0=0$) are two different scalars of two
different algebraic structures on the same variety; confusing them
produces the ``additive vs multiplicative'' paradox. The correct
statement is that both structures obey their own Grothendieck--Riemann--Roch,
each multiplicative (resp.\ additive) in its own category.
\end{observation}

\section{Six-dimensional holomorphic Chern--Simons avatar}
\label{sec:6d-hCS}

The scalar $\kappa_{\mathrm{ch}}$ has a chiral-algebraic avatar on the
six-dimensional holomorphic Chern--Simons (hCS) side. Let $Z=X\times C$
be the total space of the fibration of $X$ over a smooth curve $C$,
where the $E_1$-chiral algebra on $C$ is the boundary theory of hCS on
the bulk $6$-fold $Z$.

The Costello hCS action with target a Lie algebra $\mathfrak{g}$ is
\[
   S_{\mathrm{hCS}}[A]\;=\;\int_{Z}\Omega_{Z}\wedge
   \mathrm{tr}\!\left(\tfrac{1}{2}A\,\bar\partial A+\tfrac{1}{3}A\wedge A\wedge A\right),
\]
where $\Omega_{Z}=\Omega_{X}\wedge dz$ is the holomorphic $(4,0)$-form
on $Z$ (for $d=3$) and $A$ is a $(0,1)$-form on $Z$ with values in
$\mathfrak{g}$. BV--BRST quantisation attaches a ghost field $c$ of
ghost number $1$ and an antifield $A^{*}$, extending the classical
action to a Batalin--Vilkovisky master action $S_{\mathrm{BV}}$
satisfying $\{S_{\mathrm{BV}},S_{\mathrm{BV}}\}=2\hbar\Delta S_{\mathrm{BV}}$
on the full field--antifield space $(\mathcal{B}_{Z},Q,\Delta)$ with
$Q^{2}=0$.

\begin{proposition}[Boundary $E_{1}$-chiral algebra]
\label{prop:boundary-E1}
The boundary factorisation algebra of Costello hCS on $Z=X\times\mathbb{D}$
($\mathbb{D}$ the formal disc) is the $E_{1}$-chiral algebra
$\cA_{X}=\Phi_{d}(D^{b}(\Coh(X)))$. This is Costello's comparison
theorem (Costello 2013; Costello--Gwilliam 2021) for hCS with target
$D^{b}(\Coh(X))$ on a CY$_{d}$.
\end{proposition}

\begin{proof}
Costello's comparison identifies the boundary factorisation algebra of
hCS on $X\times\mathbb{D}$ with $\Phi_{d}(D^{b}(\Coh(X)))$ at the level
of expectation values; morphism-preservation is the chain-level
Costello--Gwilliam construction on the formal disc.
\end{proof}

\begin{proposition}[BV--BRST trace realises the Hodge supertrace]
\label{prop:bv-trace}
The one-loop effective action $W_{1}$ of Costello hCS on
$Z=X\times\mathbb{D}$ has a unique $(0,0)$-sector coefficient
$\kappa_{\mathrm{ch}}^{\mathrm{hCS}}$, read off the coefficient of the
stress-tensor operator in the boundary $E_{1}$-chiral algebra on
$\mathbb{D}$. This coefficient coincides with the Hodge supertrace,
\[
   \kappa_{\mathrm{ch}}^{\mathrm{hCS}}\;=\;\str_{\Omega^{0,\bullet}(X)}\,
   \mathbf{1}\;=\;\sum_{q=0}^{d}(-1)^{q}h^{0,q}(X)\;=\;\Xi(X)\;=\;\kappa_{\mathrm{ch}}(\cA_{X}).
\]
\end{proposition}

\begin{proof}
By Proposition~\ref{prop:boundary-E1} the boundary factorisation algebra
on $\mathbb{D}$ is $\cA_{X}$; the coefficient of the boundary
stress-tensor is $\kappa_{\mathrm{ch}}(\cA_{X})$. The one-loop BV--BRST
trace on the $(0,\bullet)$ ghost sector of $Z$ projects onto the
Dolbeault complex of $X$ in the bulk; the fermionic sign is the
cohomological $(-1)^{q}$ grading on $\Omega^{0,q}(X)$. Higher Hodge
columns $\Omega^{p,q}$ with $p\ge 1$ enter higher shadow coefficients
$S_{3},S_{4},\dots$ but not $\kappa_{\mathrm{ch}}=S_{2}$. The
identification is exactly Theorem~\ref{thm:supertrace} transposed to
the hCS ghost sector.
\end{proof}

\begin{proposition}[Tree-level Feynman realisation]
\label{prop:feynman}
The tree-level two-point Feynman amplitude of Costello hCS on $Z$,
integrated against $\Omega_{Z}$, produces the Maurer--Cartan bilinear
$[b\mid b]$-coefficient of the chiral bar complex
$B^{\mathrm{ord}}(\cA_{X})$, whose value is $\kappa_{\mathrm{ch}}(\cA_{X})$.
The Feynman propagator
$G_{\Omega}(x,y)=\Omega_{Z}^{-1}\cdot\bar\partial^{-1}\delta(x-y)$
carries the $(0,\bullet)$ Hodge column through the ghost loop and
projects onto $F^{0}/F^{-1}$ exactly as required by
Theorem~\ref{thm:supertrace}.
\end{proposition}

\begin{proof}
Tree-level Wick contraction on $Z$ with the Costello propagator,
restricted to the boundary $E_{1}$-algebra on $\mathbb{D}\subset Z$,
reproduces the chiral OPE of $\cA_{X}$. The $(0,0)$-sector of the
$[b\mid b]$ normal-ordered pairing is the Hodge supertrace on the
ghost Dolbeault complex of $X$; this is the Costello--Gwilliam
Feynman-diagrammatic identification.
\end{proof}

\begin{proposition}[$E_{3}$-avatar structure on compactly supported
Dolbeault]
\label{prop:e3-avatar}
Compactly supported Dolbeault cohomology $H^{\bullet}_{\mathrm{Dol,c}}(X)$
of a CY$_{d}$-fold $X$ carries an $E_{3}$-algebra structure for $d=3$,
realised by the hCS factorisation algebra on the bulk $3$-fold $X$
stripped of the curve direction. The Koszul averaging
\[
   \mathrm{av}\colon E_{3}\text{-}\mathrm{alg}(H^{\bullet}_{\mathrm{Dol,c}}(X))
   \;\xrightarrow{\;\int_{\Sigma_{2}}\;}\;E_{1}\text{-}\mathrm{alg}(\cA_{X})
\]
integrates the $E_{3}$-structure along the factorisation membrane
$\Sigma_{2}=\partial(B^{3})$ to produce the boundary
$E_{1}$-chiral algebra on $C$. At the level of $\kappa_{\mathrm{ch}}=S_{2}$
this averaging sends the $E_{3}$-trace on $\Omega^{0,\bullet}(X)$ to the
$E_{1}$-stress-tensor coefficient, reproducing the Hodge supertrace.
\end{proposition}

\begin{proof}
Lurie's additivity theorem identifies $E_{3}$-algebras in
chain complexes with $E_{1}$-algebras in $E_{2}$-algebras, and an
$E_{2}$-algebra in $E_{1}$-algebras is an $E_{1}$-algebra with an
additional $S^{1}$-action. On the hCS side, the $E_{3}$-structure on
$H^{\bullet}_{\mathrm{Dol,c}}(X)$ is the factorisation algebra of
$3$-dimensional holomorphic topological QFT on $X$; factorising over
$\Sigma_{2}$ (a closed surface realising the Koszul averaging membrane)
produces the $E_{1}$-algebra on $C$. The averaging map on $S_{2}$ is
an integration over $\Sigma_{2}$ that commutes with the Serre
involution $\sigma$ on the ghost sector; the character $(-1)^{d}$ of
$\sigma$ is preserved on the $E_{1}$-output scalar.
\end{proof}

\begin{proposition}[Odd-$d$ vanishing as fermionic Serre symmetry]
\label{prop:odd-d-ghost}
At odd $d$, the BV--BRST ghost complex on $X$ carries the Serre
involution $\sigma\colon\Omega^{0,q}(X)\to\Omega^{0,d-q}(X)$ as a
fermionic $\Z/2$-symmetry with character $(-1)^{d}=-1$. The one-loop
BV--BRST trace transforms as $\str\mapsto-\str$ and therefore vanishes.
Consequently
$\kappa_{\mathrm{ch}}^{\mathrm{hCS}}=0$ at every odd $d$, matching
stratum (I) of Theorem~\ref{thm:tri-stratum}.
\end{proposition}

\begin{proof}
The CY volume form induces a fibrewise pairing
$\Omega^{0,q}(X)\wedge\Omega^{0,d-q}(X)\to\Omega^{d,d}(X)\simeq\C$ on
cohomology; this pairing is the source of the Serre involution $\sigma$.
On the ghost Dolbeault complex of $X$ as a BV--BRST complex, $\sigma$
acts as a graded automorphism with parity $d\bmod 2$: it exchanges
ghost number $q$ with ghost number $d-q$. Thus $\sigma$ acts on the
one-loop trace $\str\mathbf{1}$ by the scalar $(-1)^{d}$. At $d$ odd,
$\sigma$-invariance forces the trace to be $-\str\mathbf{1}$,
hence zero.
\end{proof}

The hCS avatar makes the root cause of the stratification manifest at
the level of the BV--BRST ghost sector: Serre duality is a fermionic
symmetry of the ghost complex, and odd-$d$ vanishing of
$\kappa_{\mathrm{ch}}$ is the vanishing of a $\sigma$-antisymmetric
trace. Proposition~\ref{prop:e3-avatar} locates the Hodge supertrace
not only at $E_{1}$-chiral output but also at $E_{3}$-bulk structure on
$H^{\bullet}_{\mathrm{Dol,c}}(X)$, with the Koszul averaging map
intertwining them. The ``additivity under direct sums'' reading of the
Heisenberg-level scalar is an orthogonality of $E_{3}$-blocks on the
ghost complex; the K\"unneth-multiplicative reading of the Hodge
supertrace is the graded tensor product of ghost Dolbeault complexes.
Both obey their own Grothendieck--Riemann--Roch, each correct in its
own lane.

\section{Cross-consistency}
\label{sec:cross-consistency}

The identifications are internally consistent across the Calabi--Yau
landscape and with the $K3\times E$ scalar spectrum.

\begin{itemize}
\item \emph{K3.} Stratum (II)$\cap$(III) with $n=1$:
      $\Xi(K3)=2=\chi(\cO_{K3})$. Heisenberg level:
      $\kappa_{\mathrm{ch}}^{\mathrm{Heis}}(K3)=2$ matches the Hodge
      supertrace via rank-$2$ lattice VOA on $(1,0,1)$.
\item \emph{Elliptic curve.} $h^{1,0}(E)=1$ violates the
      stratum-hypothesis; $\Xi(E)=0$ by pairwise cancellation of
      $(1,1)$. The Heisenberg-level scalar
      $\kappa_{\mathrm{ch}}^{\mathrm{Heis}}(E)=1$ is the rank-$1$ lattice
      VOA on the generator, a different invariant.
\item \emph{$K3\times E$.} On the Hodge-supertrace lane,
      $\Xi(K3\times E)=2\cdot 0=0$ by K\"unneth. In the variety lane,
      $\chi(\cO_{K3\times E})=2\cdot 0=0$ by K\"unneth. In the
      Heisenberg-level lane,
      $\kappa_{\mathrm{ch}}^{\mathrm{Heis}}(K3\times E)=2+1=3$ by
      additive generator counting. In the Borcherds-weight lane,
      $\kappa_{\mathrm{BKM}}(\Delta_{5})=c_{1}(0)/2=10/2=5$ via the
      Gritsenko--Nikulin level-$1$ Borcherds lift. In the fibre-rank
      lane, $\kappa_{\mathrm{fiber}}(K3\times E)=\rk(\Lambda_{K3})+2=24$
      via the Mukai lattice. The five scalars $\{0,3,5,24\}$ attach to
      four distinct constructions and one variety invariant on the
      single CY$_{3}$ total space.
\item \emph{Quintic.} Hodge-supertrace: $\Xi(X_{5})=0$ by stratum (I).
      BCOV: $\kappa_{\mathrm{ch}}^{\mathrm{BCOV}}=-25/3$.
      $\chi(\cO_{X_{5}})=0$ by stratum (I). The three coexist.
\item \emph{Local $\mathbb{P}^{2}$.} $\kappa_{\mathrm{ch}}=3/2$ by
      Proposition~\ref{prop:local-p2}; no compact $\chi(\cO_X)$ matches.
\item \emph{Resolved conifold.} $\kappa_{\mathrm{ch}}=1$ by
      Proposition~\ref{prop:conifold}; assigned via compact-generator
      rank, not any Hodge column.
\item \emph{Universal Borcherds identity.} $\kappa_{\mathrm{BKM}}(\Phi_N)=c_N(0)/2$
      for $N\in\{1,2,3,4,6\}$ is the Borcherds weight theorem (Borcherds
      1995; Gritsenko 1999); it is a specialisation-stage invariant and
      is not additively related to $\kappa_{\mathrm{ch}}+\chi(\cO_{\mathrm{fiber}})$
      in any case.
\end{itemize}

\section{Open questions}
\label{sec:open}

\begin{enumerate}
\item Is the identification
      $\kappa_{\mathrm{ch}}(\Tot(\omega_{S}))=\tfrac12\chi_{\mathrm{top}}(S)$
      of Theorem~\ref{thm:stratified} universal across Fano surfaces
      $S$ (beyond $\mathbb{P}^{2}$), or are there Fano surfaces on which
      the equality fails through a non-trivial McKay-quiver
      automorphism group? The Hirzebruch surfaces $\mathbb{F}_{n}$ and
      del Pezzo surfaces $dP_{k}$ are the first test cases.
\item Is there a compact CY$_{3}$ for which the Heisenberg-level scalar
      $\kappa_{\mathrm{ch}}^{\mathrm{Heis}}$ is computable first-principles
      without K3-fibre or toric input? The compact quintic has
      $\kappa_{\mathrm{ch}}^{\mathrm{BCOV}}=-25/3$ but
      $\kappa_{\mathrm{ch}}^{\mathrm{Heis}}$ is undefined by the present
      constructions.
\item At odd $d\ge 7$, the Serre-involution argument applies uniformly;
      is there an instanton-corrected scalar in the hCS avatar, beyond
      the BCOV one-loop at $d=3$, that could lift the universal
      vanishing at $d=7,9,\dots$?
\item The conifold value $\kappa_{\mathrm{ch}}=1$ is the minimal
      nontrivial seed in the Kontsevich--Soibelman CoHA
      wall-crossing formalism. Does every non-compact CY$_{3}$ with a
      single compact $\mathbb{P}^{1}$-cycle produce
      $\kappa_{\mathrm{ch}}=1$, independently of the normal bundle
      splitting type $(\cO(-1)^{\oplus 2}, \cO(0)\oplus\cO(-2), \cO(1)\oplus\cO(-3),\dots)$?
\item Does the $E_{3}$-algebra structure on
      $H^{\bullet}_{\mathrm{Dol,c}}(X)$ of Proposition~\ref{prop:e3-avatar}
      descend from a fully extended $3$-dimensional framed TFT, and is
      its Koszul averaging to the boundary $E_{1}$ the CY-to-chiral
      functor $\Phi_{3}$ in $(\infty,1)$-categorical form? The
      candidate construction is Costello's BV-factorisation functor
      composed with the Lurie cobordism hypothesis on $\mathrm{Coh}(X)$
      as a dualisable object of $\mathrm{CY}_{3}$-cat.
\end{enumerate}

\appendix

\section{Rectification log}
\label{sec:rect-log}

\begin{itemize}
\item The phrase ``naive CY-D formula'' is replaced throughout by
      direct mathematical content: the identity
      $\kappa_{\mathrm{ch}}=\chi(\cO_{X})$ holds as a structural equality
      at K3 alone; at every other odd-$d$ compact case it is a tautology
      of Serre cancellation; at $K3^{[n]}$ ($n\ge 2$) it holds by a
      different mechanism (middle-column survival) on each side.
\item The root-cause formulation ``$\kappa_{\mathrm{ch}}$ additive under
      direct sums, not fibre products'' is corrected in
      Observation~\ref{obs:root-cause}. The Hodge supertrace is
      K\"unneth-multiplicative (not additive) on compact products; the
      appearance of an additive rule in the $K3\times E$ landscape
      comes from conflating the Hodge supertrace with the Heisenberg-level
      generator count $\kappa_{\mathrm{ch}}^{\mathrm{Heis}}$, which is
      additive in its own category. The true structural root cause is
      the character $(-1)^{d}$ of the Serre involution $\sigma$ on
      alternating sums (Lemma~\ref{lem:serre-sign}).
\item The $K3\times E$ disambiguation clarifies that $\kappa_{\mathrm{cat}}=0$,
      $\kappa_{\mathrm{ch}}^{\mathrm{Hodge}}=0$,
      $\kappa_{\mathrm{ch}}^{\mathrm{Heis}}=3$,
      $\kappa_{\mathrm{BKM}}=5$, and $\kappa_{\mathrm{fiber}}=24$ are five
      distinct invariants attached to four distinct constructions and
      one variety invariant on one CY$_{3}$.
\item The hCS/BV--BRST avatar makes the odd-$d$ vanishing a manifest
      fermionic symmetry of the ghost sector, removing any appeal to
      incidental Hodge-number cancellations.
\item The $E_{3}$-avatar Proposition~\ref{prop:e3-avatar} interweaves
      the Calabi--Yau--Fukaya--Ginzburg (CFG) $E_{3}$-structure on
      compactly supported Dolbeault with the boundary $E_{1}$-chiral
      algebra via Lurie additivity and Costello factorisation on the
      $6$-fold $Z=X\times C$.
\item The stratified formula of Theorem~\ref{thm:stratified} separates
      compact, local-surface, and conifold cases by their native
      supertrace: on $\HC^{-}_{d}$, on compactly supported Dolbeault,
      and on Kronecker--$A_{1}$ Hochschild respectively.
\item Five attack--heal cycles converged on the values
      $\{0,2,3,0,3/2,1,-25/3\}$ for
      $\{E,K3,K3^{[2]},K3\times E,\text{local }\mathbb{P}^{2},\text{conifold},\text{quintic (BCOV)}\}$,
      each verified from a native supertrace construction. Each cycle
      tested an attacker's hypothesis (Hodge-column-only root cause,
      additive--vs--multiplicative distinction, $\chi(\cO_X)$-identity
      universality, $E_{3}$-avatar rejection, six-dimensional hCS
      ghost-sign anomaly) and healed by identifying the surviving
      structural content.
\end{itemize}

\section{Primary literature}
\label{sec:primary}

\begin{itemize}
\item Beauville, A., \emph{Vari\'et\'es K\"ahleriennes dont la premi\`ere
      classe de Chern est nulle}, J.\ Differential Geom.\ \textbf{18}
      (1983), 755--782.
\item Bershadsky, M., Cecotti, S., Ooguri, H., Vafa, C.,
      \emph{Kodaira--Spencer theory of gravity and exact results for
      quantum string amplitudes}, Comm.\ Math.\ Phys.\ \textbf{165}
      (1994), 311--427.
\item Borcherds, R.\ E., \emph{Automorphic forms on $O_{s+2,2}(\mathbb{R})$
      and infinite products}, Invent.\ Math.\ \textbf{120} (1995), 161--213.
\item Borcherds, R.\ E., \emph{Automorphic forms with singularities on
      Grassmannians}, Invent.\ Math.\ \textbf{132} (1998), 491--562.
\item Candelas, P., de la Ossa, X., Green, P., Parkes, L., \emph{A pair of
      Calabi--Yau manifolds as an exactly soluble superconformal theory},
      Nuclear Phys.\ B \textbf{359} (1991), 21--74.
\item Costello, K., \emph{Supersymmetric gauge theory and the Yangian},
      arXiv:1303.2632 (2013).
\item Costello, K., Gwilliam, O., \emph{Factorization algebras in quantum
      field theory}, Cambridge Univ.\ Press, Vols.\ 1 (2017) and 2 (2021).
\item Gritsenko, V., \emph{Reflective modular forms in algebraic geometry},
      arXiv:1005.3753 (1999/2010).
\item Klebanov, I.\ R., Witten, E., \emph{Superconformal field theory on
      threebranes at a Calabi--Yau singularity}, Nuclear Phys.\ B
      \textbf{536} (1998), 199--218.
\item Klemm, A., Pandharipande, R., \emph{Enumerative geometry of
      Calabi--Yau $4$-folds}, Comm.\ Math.\ Phys.\ \textbf{281} (2008),
      621--653.
\item Kontsevich, M., Soibelman, Y., \emph{Cohomological Hall algebra,
      exponential Hodge structures and motivic Donaldson--Thomas
      invariants}, Commun.\ Number Theory Phys.\ \textbf{5} (2011), 231--352.
\item Lurie, J., \emph{Higher algebra}, online manuscript (2017),
      Theorem 5.1.2.2 (additivity).
\item Maulik, D., Okounkov, A., \emph{Quantum groups and quantum
      cohomology}, Ast\'erisque \textbf{408} (2019).
\item Nagao, K., \emph{Derived categories of small toric Calabi--Yau
      $3$-folds and counting invariants}, Q.\ J.\ Math.\ \textbf{63}
      (2012), 965--1007.
\item Nagao, K., Nakajima, H., \emph{Counting invariant of perverse
      coherent sheaves and its wall-crossing}, Int.\ Math.\ Res.\ Not.\
      \textbf{2011}, 3885--3938.
\item Ooguri, H., Yamazaki, M., \emph{Crystal melting and toric
      Calabi--Yau manifolds}, Comm.\ Math.\ Phys.\ \textbf{292} (2009), 179--199.
\item Schiffmann, O., Vasserot, E., \emph{Cherednik algebras, $W$-algebras
      and the equivariant cohomology of the moduli space of instantons
      on $\mathbb{A}^{2}$}, Publ.\ Math.\ IHES \textbf{118} (2013), 213--342.
\item Szendr\H{o}i, B., \emph{Non-commutative Donaldson--Thomas theory and
      the conifold}, Geom.\ Topol.\ \textbf{12} (2008), 1171--1202.
\end{itemize}

\end{document}
